
{\actuality} Актуальность диссертации заключается в ее междисциплинарном характере, демонстрирующем сочетание из следующих подходов: 

\begin{enumerate}[label=\fbox{\arabic*}]
	\item систематический анализ геопространственных данных из множества источников; 
	\item статистическое моделирование и обработка данных библиотеками Python и R, AWK и Octave; 
	\item геологический анализ литературы; 
	\item картографическое отображение и моделирование желобов скриптами Generic Mapping Tools (GMT) и плагинами QGIS. 
\end{enumerate}
Диссертация вносит вклад в исследования по геотектонике, геологии и геоморфологии дна Тихого океана. 

Формирование глубоководного желоба представляет собой сложный процесс, на который влияет совокупность различных геологических и тектонических факторов, объясняющих процесс его формирования и дальнейшее развитие в зависимости от степени влияния данных факторов, влияющих на морфоструктуру. Попытки исследований желобов по поперечному сечению профилей проводились в прошлом, например, но все же основанные на технических приемах того времени. На настоящий момент в существующей литературе нет работ, основанных на систематическом применения скриптов GMT для моделирования морфоструктуры океанских желобов, и это исследование способствовало этому вопросу.

В настоящее время изучение морских геологических явлений и сложных процессов с помощью программирования и скриптов является эффективным методом. Быстрое развитие передовых методов анализа данных представляет такие инструменты, как GMT, Octave/MATLAB, R и Python. Использование скриптовых методов особенно эффективно применительно к огромному количеству морских геологических данных (т.н. big data). Обработка больших массивов данных с помощью скриптовой техники картографирования и графической визуализации представляется критически важным подходом, поскольку алгоритмы библиотек предоставляют доступ к точному и быстрому анализу больших массивов данных. Конкретная информация о глубинах труднодоступных океанических желобов может быть получена для точной визуализации и анализа их подводной геоморфологии от локального до регионального и глобального масштабов.

Однако, несмотря на существующие отдельные публикации, в настоящий момент отсутствует единообразие в изучении глубоководных желобов, образовавшихся в определенных геологических и тектонических условиях на окраинах Тихого океана, а также имеет место нехватка систематического картографирования желобов применительно к Тихоокеанскому региону. Отсутствие системного подхода к картографированию глубоководных желобов приводит к недостаточному пониманию геоморфологических различий их профилей в разных частях океана: южном и северном, восточном и западном, а также и их особенностей, вызванных геологическими и тектоническими локальные условиями в местах формирования. 

В связи с этим, данная диссертация направлена на восполнение существующих на данный момент методологических пробелов путем систематического картографирования, геоморфологическое моделирования, пространственного анализа и классификации 20 желобов Тихого океана. В частности, данная работа представляет собой междисциплинарный подход, применяя сочетание следующих дисциплин для решения отдельных задач: 

\begin{enumerate}[label=(\roman*)]
	\item геодинамика и тектоника (анализ происхождения, эволюции и формирования желобов); 
	\item картография (скриптовые алгоритмы обработки пространственных данных в GMT и плагины QGIS); 
	\item статиcтическая обработка данных (библиотеки Python, R и Octave); 
	\item геоморфология (анализ современных морфоструктур рельефа дна Тихого океана). 
\end{enumerate}

В диссертации представлен системный подход к картографированию, моделированию и сравнительному анализу морфоструктур глубоководных желобов на основе сочетания передовых технических методов визуализации и анализа данных в области картографии и статистики и теоретических знаний по геодинамике и эволюции глубоководных желобов Тихого океана.

% {\progress}
% Этот раздел должен быть отдельным структурным элементом по
% ГОСТ, но он, как правило, включается в описание актуальности
% темы. Нужен он отдельным структурынм элемементом или нет ---
% смотрите другие диссертации вашего совета, скорее всего не нужен.

{\aim} данной работы является сравнительный анализ вариативности геоморфологии 20 глубоководных желобов Тихого океана в их различных сегментах на основе прикладного статистического анализа, поддерживаемого GMT, Octave, Библиотеки Python и R, и геоморфологическое моделирование желобов для классификации и ранжирования их форм рельефа на основе общих характеристик в ортогональном плане:
\begin{enumerate}
    \item Проанализировать вариации геоморфологической формы подводной геоморфологии в 20 глубоководных районах желоба Тихого океана. 
    \item Выполнить классификацию геоморфологических типов 20 глубоководных желобов на основе их геометрических свойства, отличия и сходство. 
\end{enumerate}
Следующие типы форм геоморфологических профилей желобов были выявлены и подразделены на семь типов: U-образные (в плане), V-образные (в плане), асимметричный, серповидный, извилистый, удлиненный, каскадный. Для каждого типа (U, V, асимметричный и т. д.) выделены характерные подтипы крутизны: сильная, очень сильная, экстремальная, крутая, очень крутой. Склоны долин классифицируются следующим образом: очень высокие, высокие, умеренные, низкие, в зависимости от кривизны и степени. Классы размеров и уклонов долин анализируются в контексте физической среды, тектонических и геологических условий района желоба.

Для~достижения поставленной цели необходимо было решить следующие {\tasks}:
\begin{enumerate}[label=\Roman*.]% https://tex.stackexchange.com/a/476052/104425
  \item Моделирование морфологии продольно-поперечных разрезов (профилей) глубоководных желобов с помощью модулей скриптового картографического инструмента GMT для систематического построения ряда профилей поперечного сечения для сравнительного анализа геоморфологии 20 желобов.
  \item Картографирование локальных топографических и геологических условий в районах расположения 20 глубоководных желобов, на основе имеющихся наборов данных.
  \item Картографирование геофизических свойств (морской гравитации в открытом воздухе и модели геоида) 20 желобов на основе имеющихся наборов данных. 
  \item Выполнение статистического анализа батиметрических продольно-поперечных разрезов (профилей) глубоководных желобов на основе использования методов алгоритмов программирования (Python и R) для получения моделей взаимосвязи геологических факторов и геоморфологии желобов.
  \item Методическая разработка совместного применение статистических библиотек для анализа и визуализации данных (Python, R, GMT, Octave, AWK, QGIS) для анализа данных и моделирования
  \item Анализ вариативности морфологии желобов вдоль окраин Тихого океана для оценки взаимосвязи геодинамических факторов, влияющих на их морфоструктуру. 
\end{enumerate}

{\novelty}
\begin{enumerate}[beginpenalty=10000] % https://tex.stackexchange.com/a/476052/104425
  \item Впервые разработаны экспериментальные методы GMT, примененные к моделированию именно глубоководных желобов, а также расширении существующих методов пространственного моделирования с использованием языков программирования Python, R, AWK и Octave для моделирования, статистического анализа, визуализации и геоморфологической классификации подводных форм рельефа желобов. Использование методов расширенного анализа данных оказало решающее значение для полученной точной и надежной обработки данных, поскольку понимание форм рельефа морского дна может быть основано исключительно на компьютерном моделировании данных из-за их недоступности месторасположения желобов. 

  \item Техническая инновация работы заключается в междисциплинарном подходе, сочетающем ГИС-анализ и статистические методы (R, Python, MATLAB), которые способствуют изучению океанских желобов, при пространственном анализе больших данных. Моделирование закономерностей корреляции геоморфологической структуры и геологических особенностей морского дна в районе желобов имеют решающее значение для моделирования рельефа дна океана. Выбор методологии, инструментов и алгоритмов объясняется целями и задачами исследования, поскольку специфика морской геологии, геодинамики и геофизики заключается в особо высоких требованиях к методическому уровню обработки данных, применении самых передовых технологий и инструментария для точного анализа данных, недоступных для непосредственного изучения в связи с удаленностью объекта (глубоководные желоба). 
  
\item Было выполнено оригинальное исследование на основе наборов данных, которые обрабатывались, вычислялись и анализировались в полуавтоматическом режиме с помощью алгоритмов машинного обучения в картографии. Таким образом, диссертация представляет комплексный подход. Основанные на глубинном анализе геологического строения желобов (представленный в описательной главе 2) и технические подходы, использующие современные алгоритмы анализа данных и эффективную визуализацию путем использования языков программирования и скриптовой картографической программы GMT, не применяемого ранее в существующих работах по сравнительной морфологии глубоководных желобов Тихого океана. 

\item Впервые представлена систематическая классификации и осуществлено сравнительное моделировании геоморфологических профилей глубоководных желобов Тихого океана, выполненных с помощью последовательного использования современных наборов инструментов для создания картографических скриптов и статистического анализа. Техническая новизна заключается в сочетании GIS, GMT, Python, AWK, R применительно к задачам картографирования и моделирования подводных форм рельефа дна глубоководных желобов Тихого океана.  
\end{enumerate}

{\influence} Геоморфология рельефа дна Мирового океана в целом и океанических желобов в частности (как особых его структур), их образование, формирование и эволюция – вопрос особой важности для научного сообщества в области геодинамики, тектоники, морской геофизике и геологии. Информация о геоморфологии глубоководных  желобов имеет важное значение для изучения динамической обстановки и тектонической активности окраин Тихого океана. Актуальность этого вопроса значительно увеличилось с начала бурного развития информационных технологий (ИТ), инструментов и методов расширенного анализа данных, но их понимание тем не менее остается неоднородным. Поскольку большинство желобов расположены по окраинам именно Тихого океана (по сравнению с другими), изучение тихоокеанских желобов играет центральную роль для анализа и понимания их образования. 

Особые геологические условия, наличие зон тектонической субдукции, обширная территория Тихого океана со сложной системой циркуляции (Рис. 1), распространение зоны «огненного кольца», т.е. сейсмически активного пояса землетрясений и вулканических процессов, делают желоба Тихого океана чувствительными к факторам, влияющим на их формирование, которое вызывает изменения в их геоморфологических формах. До 80\% землетрясений, зафиксированных сейсмическими службами, сосредоточены в Тихоокеанском «огненном кольце» [6, 7]. Тектонические структуры дня Тихого океана отличаются вариативностью и сложностью форм. В общей классификации можно выделить острова, подводные окраины материков, океанические валы, подводные горы, срединно-океанические хребты, котловины, каньоны, абиссальные равнины, ложе океанических впадин, подводные внутриокеанические каньоны, а также глубоководные желоба, зоны разломов, срединноокеанические хребты и срединноокеанические рифты, которые входят в единую систему рифтов растяжения. В этом контексте, наиболее репрезентативными индикаторами вариаций глубоководных желобов являются геологические и тектонические факторы, такие как динамика литосферной коры, влияющая на скорость и интенсивность субдукции плит, величину и повторяемость подводных вулканов, вызывающих активную седиментацию. 

Геоморфология морского дна в желобах зависит от множества факторов, которые в конечном счете влияют на форму их рельефа, определяют вариативность морфоструктур и динамику. Анализ геоморфологии желобов имеет практическое значение при изучении тектонически активных зон Тихого океана. Эта диссертация рассматривает такие факторы как формирование тектонических плит, субдукцию, историческое геологическое развитие, расположение основных зон землетрясений и подводных вулканов как основные типы ключевых факторов, влияющих на формирование океанических желобов. Вторичные факторы включают океанические течения, седиментацию и особенности распределение биоты, способствующие седиментации. 

Знание и правильное понимание движущих факторов, воздействующих на геоморфологию океанического дна, дает понимание геоморфологии океанского дна. Учитывая важность глубоководных районов, изучение морфоструктур желобов является актуальной задачей. Используя методы моделирования данных, были сопоставлены и проанализированы формы пересечения профилей 20 желобов в ортогональном направлении для выделения различий и вариации в их геоморфологии. 

{\methods} Мультидисциплинарная интеграция картографических методов исследования (GIS, GMT), статистической обработки данных (Python, AWK, R) и анализа обширного материала литературных источников по геологии и тектонике Тихого Океана применительно к задачам картографирования и моделирования подводных форм рельефа дна глубоководных желобов Тихого океана.

{\defpositions}
\begin{enumerate}[label=\emph{\alph*})]
 	 \item Первое положение. Научное обоснование, комплекс картографических методических решений и скриптовых алгоритмов GMT для моделирования морфоструктур глубоководных желобов Тихого океана, находящихся под воздействием региональных тектонических условий и меняющихся скоростей движения литосферных плит.
  	\item Второе положение. Основные пути повышения: i) методической актуальности (численные методы, скриптовые языки и программирование на языках Python, R, Octave); ii)  наукоемкости (анализ обширного объема материалов как картографического, так и текстуального типа); iii) информативности (обработка больших массивов геолого-геофизических и геоморфологических данных с неочевидными, нелинейными связями, big data) для комплексных исследований дна океана при совокупном влиянии тектонических и геодинамических факторов, определяющих морфоструктуру дна Тихого океана.
 	 \item Третье положение. Теоретическая аргументация методологических решений по созданию методики комплексного картографирования профилей желобов двух- и трехмерным моделированием (2D и 3D-моделирование), описывающих: а) изменение крутизны и уклона профилей при вариативности геологических и тектонических условий; б) изменение геофизических свойств области исследования при вариациях гравиметрии вдоль русла желобов.
 	 \item Четвертое положение. Современный геоинформационных методический комплекс с применением скриптовых решений GMT и статистического анализа на языках Python, R и Octave для исследований изменения морфологических, геометрических и геофизических свойств желобов под воздействием геолого-геофизических региональных эффектов от воздействии вариативности скорости движения и субдукции литосферных плит.
  	 \item Пятое положение. Западный и восточный регионы Тихого океана характеризуются вариативностью типов глубоководных желобов и маркированы U, V-образными и асимметричными формами разрезов поперечного сечения. Сейсмически активные зоны и геолого-геофизические условия определяют морфологию желобов.
\end{enumerate}

%{\reliability} полученных результатов обеспечивается \ldots \ Результаты находятся в соответствии с результатами, полученными другими авторами.

{\probation}
По результатам исследований, выполненных в рамках данной диссертации, опубликовано 70 основных печатных работ, из них 21 статей в периодических изданиях, индексированных в Scopus и Web of Science и входящих в список рекомендованных ВАК, 3 тезиса докладов материалов конференций и 46 публикаций в прочих журналах.

{\contribution} Автор принимала активное непосредственное участие в работе.
\begin{enumerate}[label=\emph{\Alph*})]
	\item Автор лично написала скрипты GMT, использованные для моделирования желобов и тематического картографирования, а также скрипты на языках программирования Python, R и Octave для дополнительного статистического анализа батиметрии желобов. С использованием данных скриптов автором лично проведено картографирование, 2D и 3D моделирование всех основных 20 желобов, представленных в настоящей диссертации, геологическое описание исследованных районов желобов Тихого океана. 

	\item Автором лично выполнены все этапы картографирования, включающие следующие шаги: сбор, обработка, анализ, интерпретация и визуализация данных и моделирования желобов. Автором лично выполнен отбор и предварительный анализ данных на пригодность (разрешение, надежность источника, актуальность) из доступных геофизических, геологических и топографических источников данных. Автором лично проведена обработка данных для картографирования (форматирование, проецирование и подготовка отдельных слоев по координатам профилей) и последующее систематическое картографирование. 

	\item Все иллюстрации без исключения (как карты, так и графики), представленные в диссертации, являются результатом личной работы автора: картографирование на основе GMT и частично QGIS с использованием плагинов для цифрования, графики профилей желобов, оцифрованные в GMT по авторской разработанной методике с использованием модулей GMT, статистические графики, выполненные на основе R, Python и Octave. При осуществлении геоморфологического моделирования использовалась авторская разработка профилирования желобов с помощью комбинации модулей GMT, позволяющая с высокой скоростью и большой точностью производить оцифровку профилей по параметрам, разработанным автором, представленным в данной диссертации и апробированным на 20 желобах Тихого океана.

	\item По выполнении картографирования и технических скриптовых работ, автор лично провела интерпретацию морфологии желобов на основе результатов моделирования, а также анализ вариативности геофизических и геологических условий в областях исследования (мест расположения желобов) на основе визуализации данных. В соответствии с целью и задачами диссертации, автор проанализировала и обобщила обширный фактический материал (как картографические, так и литературные источники), опубликованные в предыдущих существующих работах. Результаты графического моделирования и картографирования были использованы автором как иллюстрации в публикациях. 
\end{enumerate}

\ifnumequal{\value{bibliosel}}{0}
{%%% Встроенная реализация с загрузкой файла через движок bibtex8. (При желании, внутри можно использовать обычные ссылки, наподобие `\cite{vakbib1,vakbib2}`).
    {\publications} Основные результаты по теме диссертации изложены
    в 70 печатных изданиях, 21 из которых изданы в журналах, рекомендованных ВАК, Scopus и Web of Science,
    3 тезиса докладов материалов конференций и 46 публикаций в прочих журналах.
}%

